\documentclass[11pt,letterpaper]{article}
\usepackage[margin=0.85in,headheight=16pt]{geometry}
\usepackage[T1]{fontenc}
\usepackage[utf8]{inputenc}
\usepackage{lmodern,microtype}
\usepackage{amsmath,amssymb,mathtools,booktabs,longtable,tabularx,array}
\usepackage[dvipsnames]{xcolor}
\usepackage{hyperref,fancyhdr,titlesec,enumitem,graphicx,tcolorbox}
\usepackage{lastpage}
\definecolor{ink}{HTML}{14263B}
\definecolor{sea}{HTML}{167C8A}
\definecolor{gold}{HTML}{C49747}
\definecolor{pale}{HTML}{EEF5F6}
\definecolor{graytext}{HTML}{516275}
\hypersetup{colorlinks=true,linkcolor=ink,urlcolor=sea,bookmarksopen=true,pdfauthor={Independent Educational Engineering Suite}}
\titleformat{\section}{\Large\bfseries\color{ink}}{\thesection}{0.7em}{}[\vspace{-0.3em}\color{gold}\titlerule]
\titleformat{\subsection}{\large\bfseries\color{sea}}{\thesubsection}{0.6em}{}
\setlist[itemize]{leftmargin=1.3em,itemsep=2pt,topsep=3pt}
\setlist[enumerate]{leftmargin=1.5em,itemsep=2pt,topsep=3pt}
\pagestyle{fancy}
\fancyhf{}
\fancyhead[L]{\small\sffamily\color{graytext}SECURITY ANALYSIS / ENGINEERING SUITE}
\fancyhead[R]{\small\sffamily\color{graytext}\textbf{\thepage}/\pageref{LastPage}}
\renewcommand{\headrulewidth}{0.3pt}
\renewcommand{\footrulewidth}{0pt}
\newtcolorbox{methodbox}[1][]{colback=pale,colframe=sea,boxrule=0.55pt,arc=1mm,title={\bfseries Method note},fonttitle=\sffamily,#1}
\newtcolorbox{sourcebox}[1][]{colback=white,colframe=gold,boxrule=0.55pt,arc=1mm,title={\bfseries Source boundary},fonttitle=\sffamily,#1}
\newcommand{\term}[1]{\textbf{\color{ink}#1}}
\newcommand{\src}{\textsf{\color{sea}[SOURCE]}}
\newcommand{\ext}{\textsf{\color{gold}[EXTENSION]}}
\newcommand{\gap}{\textsf{\color{red}[SOURCE GAP]}}
\newcommand{\answer}[1]{\par\smallskip\noindent\textbf{Check:} #1\par}
\newcommand{\cover}[3]{%
\begin{titlepage}\pagecolor{ink}\color{white}\vspace*{1.15in}
{\sffamily\large\color{gold}INDEPENDENT TECHNICAL DOCUMENTATION SUITE\par}\vspace{0.4in}
{\Huge\sffamily\bfseries #1\par}\vspace{0.15in}
{\Large\sffamily #2\par}\vspace{0.65in}
\color{gold}\rule{\textwidth}{2pt}\par\vspace{0.4in}
{\large #3\par}\vfill
{\small\sffamily Based on a critical reading of Graham and Dodd's \emph{Security Analysis}, Seventh Edition (2023). Original educational analysis; not affiliated with the publisher.\par}
\end{titlepage}\nopagecolor\color{black}}
\newcommand{\front}[3]{\cover{#1}{#2}{#3}\tableofcontents\clearpage
\begin{sourcebox}\src\ means a concept is attributable to the relevant chapter or seventh-edition contributor, summarized independently rather than reproduced. \ext\ means a modern method or engineering recommendation prepared for this suite, not an original Graham--Dodd rule. Historical examples and standards should not be treated as current market requirements. All numerical examples are synthetic, educational, and not securities recommendations.\end{sourcebox}}
\setlength{\parskip}{5pt}\setlength{\parindent}{0pt}
\begin{document}
\front{Financial Modeling}{Applied Exercises and Verified Answer Keys}{Twelve fictional cases in earnings quality, cash flow, bonds, recovery and portfolio liquidity}

\section{Workbook instructions}
All cases use fictional issuers and amounts. Unless specified, currency is USD and amounts are in millions. Present (1) inputs; (2) formulas; (3) units; (4) base/bear alternatives; (5) investment decision and uncertainty. Each exercise has an independently reproducible result check later in this document.

\section{Case 1: Earnings quality and normalization}
\textbf{Data:} Reported EBIT 150; litigation settlement gain 20 included in EBIT; recurring annual restructuring costs of 12 included in EBIT; statutory tax rate 25\%; maintenance capital needs unchanged. \textbf{Tasks:} (a) calculate sustainable EBIT with no speculative cost add-back; (b) compute normalized NOPAT; (c) explain why recurring restructuring is not exceptional.

\section{Case 2: FCFF bridge and working-capital intensity}
\textbf{Data:} EBIT 180, tax 25\%, depreciation 30, total CapEx 55, increase in operating net working capital 20. \textbf{Tasks:} calculate FCFF; recalculate if inventory policy releases 10 of working capital; explain sustainability concerns.

\section{Case 3: Three-year DCF}
\textbf{Data:} FCFFs 80, 90, 100; WACC 10\%; perpetual growth 3\%; net debt 400; diluted shares 100. \textbf{Tasks:} compute terminal value, EV, equity value/share, and terminal value's percentage of enterprise value.

\section{Case 4: Reverse valuation}
\textbf{Data:} An equity instrument trades at \$40; your value/share interval from plausible forecasts is [\$32, \$46]. \textbf{Tasks:} calculate base margin of safety using \$40 as your selected base value; calculate lower-bound margin; state a no-go rule independent of upside.

\section{Case 5: Graham--Dodd-style net-current-asset screen}
\textbf{Data:} Cash 30, trade receivables 40, inventory 50, total liabilities 70, 10m common shares; share price \$3. \textbf{Tasks:} calculate NCAV and NCAV/share, margin of safety to NCAV, and caveats.

\section{Case 6: Liquidation realization}
\textbf{Data:} Use Case 5. Recover 100\% cash, 85\% receivables, 50\% inventory; liquidation costs 7, all liabilities paid at par. \textbf{Tasks:} compute per-share net recovery and compare with market price.

\section{Case 7: Plain-vanilla fixed-rate bond}
\textbf{Data:} Three annual 5\% coupons on \$1,000 face, 7\% required nominal yield, bullet maturity. \textbf{Tasks:} derive clean theoretical price under annual conventions; show reason for below-par price; compute effect of $+100$ bp yield with full revaluation.

\section{Case 8: Simplified waterfall recovery}
\textbf{Data:} Recoverable enterprise proceeds after costs 120; secured debt 70, unsecured debt 100, preferred 25, ordinary equity residual. Treat obligations as strictly tiered without cross-collateral complexity. \textbf{Tasks:} calculate dollar and percentage recovery for each class.

\section{Case 9: Convertible parity}
\textbf{Data:} Face 1,000, conversion price 40 per share, stock 48, convertible price 1,260. \textbf{Tasks:} compute conversion shares, equity parity value and parity premium; discuss credit/floor and dilution risk.

\section{Case 10: Leverage and capital structure}
\textbf{Data:} Asset-side enterprise value 500; debt 350; liquid excess cash 20; diluted shares 10m. A 20\% EV shock affects debt's recovery expectations but not its contractual par claim. \textbf{Tasks:} calculate equity per share before and after the EV decline (hold debt par constant for simplified exercise); compute percentage equity decline.

\section{Case 11: Two-state portfolio expected loss}
\textbf{Data:} Credit A EAD 8m, annual PD 3\%, LGD 50\%; Credit B EAD 12m, PD 5\%, LGD 40\%. \textbf{Tasks:} compute combined annual expected credit loss under additive, standalone EL assumptions; explain why portfolio tail risk cannot be inferred from the sum alone.

\section{Case 12: Liquidity and concentration stress}
\textbf{Data:} Portfolio NAV 100, cash 5, security A holding 25 with 2\% normal bid discount and 15\% stress discount, security B 20 with 1\% normal and 7\% stress discount. Remaining NAV 50 is assumed immediately realizable at par solely for this simplified calculation. A mandatory withdrawal of 10 occurs during the stress. \textbf{Tasks:} value liquidation proceeds, determine whether cash plus stressed A proceeds can fund withdrawal, and calculate residual after execution on 5 units of A market value.

\section{Answer key with checks}
\subsection{Cases 1--3}
\textbf{1.} Normalized EBIT $=150-20=130$; recurring restructuring expense 12 stays included. NOPAT $=130(0.75)=97.5$. Reversing restructuring would overstate maintainable results unless persuasive evidence shows elimination.

\textbf{2.} $FCFF=180(0.75)+30-55-20=90$. A working-capital release of 10 changes the \emph{single-year} figure to 100; do not capitalize a transient release as permanent annual cash generation.

\textbf{3.} Terminal value 1,471.43; EV $\approx 1{,}327.74$; equity 927.74; value/share $\approx\$9.28$. Discounted terminal value $\approx1{,}105.51$, about 83.3\% of EV, signaling heavy terminal assumption reliance.

\subsection{Cases 4--6}
\textbf{4.} At price 40 and selected base value 40, MOS=0\%. Lower-bound MOS $=1-40/32=-25\%$. A position can fail a precommitted margin or liquidity gate despite possible upside to 46.

\textbf{5.} $NCAV=30+40+50-70=50$; NCAV/share=\$5; at price \$3, MOS to NCAV $=1-3/5=40\%$. This is not a true liquidation discount because no recovery haircut or costs have been applied.

\textbf{6.} Realizable asset proceeds $=30+34+25=89$; net after liabilities and costs $=89-70-7=12$; \$1.20/share. A \$3 purchase loses 60\% versus this hypothetical liquidation realization before tax and time discount.

\subsection{Cases 7--9}
\textbf{7.} $P\approx\$947.51$; discount to par follows coupon 5\% below yield 7\%. At 8\% yield, $P=50/1.08+50/1.08^2+1050/1.08^3\approx\$922.69$; about 2.62\% price decline.

\textbf{8.} Secured receives 70/70=100\%; unsecured receives remaining 50 of 100=50\%; preferred and common receive 0. This is a pedagogical ordering, not a prediction of court allocation.

\textbf{9.} Shares $q=1000/40=25$; parity=25(48)=\$1,200; premium $=1260/1200-1=5\%$. The convertible can lose both debt and stock value under issuer distress.

\subsection{Cases 10--12}
\textbf{10.} Initial equity=$500+20-350=170$, or \$17/share; after EV falls to 400, equity=$400+20-350=70$, or \$7/share. Equity declines about 58.8\% in the deliberately simplified static-par model.

\textbf{11.} Annual EL A=$8m(0.03)(0.50)=\$0.12m$; B=$12m(0.05)(0.40)=\$0.24m$; total \$0.36m. Portfolio tail risk also depends on default correlation and simultaneous recovery compression.

\textbf{12.} Stressed proceeds A=$25(0.85)=21.25$, B=$20(0.93)=18.60$, cash 5 and other 50, total 94.85. Cash+all stressed A proceeds=26.25, sufficient to fund withdrawal. Selling 5 of pre-stress A notional at a 15\% discount yields 4.25 cash; combined with cash=9.25, leaving a 0.75 shortfall unless further assets are sold. State whether the 10 withdrawal is gross or net of fees.

\section{Laboratory quality criteria}
Each submitted spreadsheet or program must flag mismatched periods, break on missing values rather than substituting zero, preserve source figures, isolate scenarios from base inputs, reconcile EV and equity claims, and provide a test that catches a 100x percent-versus-decimal error. Include a one-page investment view stating why apparent undervaluation might \emph{not} be exploitable.

\end{document}
